% Created by Hunter G. Hodges
% X -> @HuntHod
% LinkedIn -> www.linkedin.com/in/hunter-hodges

\documentclass[a4paper,11pt]{extarticle} 

\usepackage{enumitem} % for list formatting
\usepackage{geometry} % For page margins
\usepackage{titlesec} % For section title formatting
\usepackage{fancyhdr} % for Footer formatting


\geometry{a4paper, margin=1in} % Set paper size and margins

\titleformat{\section}{\bfseries \large}{\thesection}{1em}{} % Section title format
\titlespacing*{\section}{0pt}{\baselineskip}{\baselineskip} % Section title spacing


\begin{document} % starts the document ;)

\pagestyle{fancy} %Sets page style for headers
\fancyhead{} %clears all header fields
\renewcommand{\headrulewidth}{0pt} %Gets rid of the header line
\fancyfoot{} % clear all foot fields
\fancyfoot[L]{RANK LAST NAME/ORG/OFFICE/DSN: XXX-XXXX/INITIALS/DATE} % Please add author information here


\newlist{myBullets}{itemize}{3} %Creates custom list
%\setlist[myBullets]{itemsep=0pt} % Removes item separation -- you can remove comment and change pt value to add spacing
\setlist[myBullets, 1]{label=$-$, leftmargin=\parindent, rightmargin=10pt}
\setlist[myBullets, 2]{label=$--$, leftmargin=15pt, rightmargin=15pt}
\setlist[myBullets, 3]{label=$---$, leftmargin=30pt, rightmargin=30pt}


%~~~~~~~~~ This is where you change the title of the document ~~~~~~~~~~~~~%
\begin{center}
    \textbf{\large BULLET BACKGROUND PAPER}\\[12pt]
    \textbf{\large ON}\\[12pt]
    \textbf{\large WRITING BULLET BACKGROUND PAPERS}\\  % !Only change this line of the title! %
\end{center}
\medskip

%~~~~~~~~~ This is where you change the title of the document ~~~~~~~~~~~~~%

%~~~~~~~~~ Below is the purpose section ~~~~~~~~~~~~~%

\section*{PURPOSE}

Discuss the functions, format and tips for building a Bullet Background Paper (BBP). The purpose statement informs readers on both the purpose and main points of the paper.  It may be a single sentence, as in this example or a short paragraph in length.

%~~~~~~~~~ Below is the background section ~~~~~~~~~~~~~%

\section*{BACKGROUND}

\begin{myBullets}
    \item Concise background information on a single idea
    \begin{myBullets}
        \item More details on the proceeding statement
        \begin{myBullets}
            \item Even more details of the 2nd level statement. 
            \item Use 2nd and 3rd levels sparingly
        \end{myBullets}
    \end{myBullets}
    \item Chronological of a problem
\end{myBullets}

%~~~~~~~~~ Below is the discussion section ~~~~~~~~~~~~~%

\section*{DISCUSSION}

\begin{myBullets}
    \item Flexibile to save space and/or conform to a user's needs
    \begin{myBullets}
        \item Section can be titled ``Function" or ``Format" or ``Tips" or ``Details" or ``Discussion"
        \item If using 2nd and 3rd level, make sure there are at least two bullets
        \item Such as, this example
    \end{myBullets}
    \item This should be the ``main course" of the document!
\end{myBullets}


%~~~~~~~~~ Below is the recommendations section ~~~~~~~~~~~~~%

\section*{RECOMMENDATIONS}


\begin{myBullets}
    \item This section can be titled ``Conclusion" as it ends the document
    \item Recommendations are good for moving forward but write to the intended audience
    \item Be actionable, measurable, and attainable!
    \item Try to list OPR and OCRs per recommendation

\end{myBullets}


\end{document} %try to keep this to 1 to 2 pages anymore and it is just too much! 
